% Generated draft: requires mathematical and source review.
\documentclass[11pt]{article}
\usepackage[margin=25mm]{geometry}
\usepackage{fontspec}
\setmainfont{Latin Modern Roman}
\usepackage{amsmath,amssymb,mathtools,unicode-math}
\setmathfont{Latin Modern Math}
\usepackage{enumitem,needspace,xurl,hyperref}
\setlength{\parindent}{0pt}
\setlength{\parskip}{5pt}
\setlength{\emergencystretch}{3em}
\newcommand{\N}{\mathbb{N}}
\newcommand{\Z}{\mathbb{Z}}
\newcommand{\Q}{\mathbb{Q}}
\newcommand{\R}{\mathbb{R}}
\newcommand{\C}{\mathbb{C}}
\newcommand{\F}{\mathbb{F}}
\newcommand{\A}{\mathbb{A}}
\newcommand{\PP}{\mathbb{P}}
\newcommand{\E}{\mathbb{E}}
\newcommand{\eps}{\varepsilon}
\newcommand{\dd}{\,\mathrm{d}}
\begin{document}
\section*{UnsolvedMath ID: 20000218}
\textbf{Source problem code:} catalogue entry
\subsection{Definitions and mathematical statement}
Let \(G/P\) be a flag variety, \(W\) its Weyl group and \(W_P\) the parabolic Weyl subgroup; write \(X_s\) for the Schubert variety indexed by \(s\in W/W_P\). For \(S\subset W/W_P\), put \(X_S=\bigcup_{s\in S}X_s\), with the reduced union structure. The source asks which \(S\) deform to irreducibles, and whether Cohen--Macaulayness of \(X_S\) suffices. An embedded interpretation asks for a flat family of closed subschemes of \(G/P\), over an integral pointed curve, with special fiber \(X_S\) and irreducible general fiber. The aid instead requires an integral general fiber; reducedness, geometric integrality and the base-field convention are unresolved distinctions in the source.

The damaged first test case is interpreted as \(G/P\simeq(\PP^1)^n\), following the attached aid. Fix a base point in each factor, and for \(I\subseteq\{1,\ldots,n\}\) let \(X_I\) vary in the factors indexed by \(I\) and equal the base point elsewhere. For a pure \(r\)-dimensional union, its nonredundant support \(\mathcal B\) consists of \(r\)-element subsets.

The aid reports that such a union deforms to an integral variety exactly when \(\mathcal B\) has a degree-one algebraic-matroid realization. One explicit meaning is an integral \(r\)-fold \(Z\subset(\PP^1)^n\) for which projection to the factors in an \(r\)-subset \(I\) is dominant of generic degree one precisely for \(I\in\mathcal B\), and nondominant otherwise. Equivalently its coordinate functions have these bases and generate the function field already on each basis. The aid further reports, in characteristic zero, equivalence with \(\mathcal B\) being the bases of a linearly representable matroid. The field and deformation hypotheses of this classification are not supplied in full.

Its separating example in four factors is
\[
X_{\{1,2\}}\cup X_{\{2,3\}}\cup X_{\{3,4\}}.
\]
The aid reports this reduced pure surface union to be arithmetically Cohen--Macaulay, using a saturated height-two Cox ideal and a Hilbert--Burch resolution, but not deformable to an integral surface in the intended embedded sense. Its support fails matroid basis exchange: for bases \(\{1,2\}\), \(\{3,4\}\), removing \(2\) from the first and adding either \(3\) or \(4\) produces neither a supported basis. It reports this counterexample as characteristic free and four as the least number of factors, since every pure union for \(n\leq3\) is representable and deformable. The actual Cox ideal and resolution are absent.

The source's adjacent Remark 2.3 reports, in an unidentified speaker's first person, that a degeneration of Kazhdan--Lusztig varieties is a union of affine spaces and calls this an easier related result. It does not provide a theorem statement, hypotheses or attribution sufficient to use that remark as a classification result.
\subsection{Short English statement}
Classify reduced Schubert unions that occur as limits of irreducible varieties, and test whether Cohen--Macaulayness suffices. In the reconstructed product-of-projective-lines case, the attached aid reports a degree-one matroid criterion and a four-factor Cohen--Macaulay counterexample, while the original deformation category and field remain unspecified.
\subsection{Sources}
\begin{itemize}
\item Local UnsolvedMath registry, permanent ID \(20000218\), exact public code \texttt{catalogue entry}; canonical location \texttt{aim-algebraic-geometry-notes.json}, notes[217].
\item AIM workshop ``Degenerations in algebraic geometry,'' extracted Problem 2.2 and adjacent Remark 2.3; supplied source URL: \url{https://aimath.org/pastworkshops/degenalggeomproblems.pdf}. The PDF is cited from the supplied record only; it was not consulted in this drafting step. No proposer or year is supplied; the literature field is empty.
\item Attached research attempt 1, ``Matroid classification and a sharp four-factor Cohen-Macaulay obstruction,'' referring to Cox and Hilbert--Burch calculations without complete bibliographic citations or their explicit ideal/resolution.
\end{itemize}
\end{document}
